\section{Grid Interconnect Requirements}

The load demands described above provide unprecedented interconnection challenges due to the scale and volatility of their power demand. Grid operators are starting to require greater load flexibility, controllability, and predictability to maintain power grid stability, avoid overbuilding for peak demand, and reduce stress on transmission infrastructure.

Unmitigated step changes in AI workloads, especially from large, synchronized GPU clusters, can cause rapid voltage and frequency deviations outside of acceptable bands. These events risk violating interconnection requirements, degrading grid performance, and triggering delays or denials in interconnection approvals.

To secure timely grid access and ensure stable long-term operation, AI compute loads should incorporate compute settings adjustments, active power conditioning, and grid-compliance controls over Ethernet, including:

\begin{itemize}
\item \textbf{Energy storage systems:} energy storage (longer and shorter duration) with fast, real-time power compensation to stabilize power consumption, control ramp rates, and mitigate large step changes.

\item \textbf{GPU performance tuning and workload pacing:} Adjust compute firmware and operational settings to smooth rapid power fluctuations, limit cycle-to-cycle ramp rates, and suppress peak power behavior.

\item \textbf{Coordinated control strategies:} Integrate storage, compute, and facility power distribution systems to ensure compliance with utility requirements for ramp rate, transient stability, harmonics/flicker, and voltage ride-through (VRT), while supporting overall grid stability.
\end{itemize}

Looking ahead, AI Factories have the potential to evolve into grid-supporting assets rather than operating solely as large, passive loads. By adopting advanced capabilities, industry partners can help ensure that AI infrastructure not only scales sustainably but also strengthens the energy systems it depends on.

\begin{itemize}
\item \textbf{Grid-forming and fast-response controls via Ethernet} that actively support voltage and frequency stability during grid disturbances.

\item \textbf{Enhanced fault ride-through capability for internal IT loads,} surpassing current ITIC curve expectations by referencing ERCOT Voltage Ride-Through (VRT) requirements for large electronic loads, to maintain critical operations and support grid recovery during voltage sags or faults.

\item \textbf{Real-time power support,} delivering supplemental power when the grid is under stress and absorbing excess power when available, helping smooth sudden changes and maintain utility grid stability.

\item \textbf{Resilient system designs and operational practices} that position AI Factories as contributors to overall grid reliability and resilience, enabling faster, more predictable scaling of deployments.
\end{itemize}

To ensure grid stability and power quality in high-density AI Factories, energy storage must be strategically deployed at both ends of the electrical architecture. Energy storage, such as a Battery Energy Storage System (BESS), is placed adjacent to the utility interconnection and on-prem generation to manage steady-state power fluctuations from large-scale GPU clusters. The BESS provides load averaging, grid-forming support, and transitional power during utility-to-generator switching events. At the other end, short-duration storage such as capacitors are integrated near the compute racks to mitigate fast dynamic power behavior (as fast as 400 microseconds) generated by AI workloads. These devices serve to limit ramp rates and absorb power peaks and troughs, ensuring a stable and predictable 800 VDC supply to sensitive compute infrastructure. This dual-layered storage strategy is critical to meeting the performance and reliability demands of AI-native power systems.

Interconnection approvals remain a major bottleneck for AI Factory deployment, largely due to the unpredictable nature of GPU-driven load swings. Industry collaboration is needed to establish standardized load behavior profiles, response metrics, and conditioning requirements. Clear alignment on ramp-rate limits, load flexibility, and energy storage integration will help utilities evaluate AI Factory projects more confidently, speeding approvals and enabling faster, more scalable deployment.

Energy storage must be treated as a core element of AI Factory architecture. Beyond stabilizing GPU loads, its standardized roles should cover transient mitigation, backup source, and load shaping, with common integration interfaces for seamless deployment at rack, row, or facility level. This approach ensures interoperability with power conversion systems and grid interconnection schemes, improves equipment utilization, and enables reliable scaling without compromising compute performance.
